\documentclass{beamer}
\mode<presentation>
\usepackage{amsmath,amssymb,mathtools}
\usepackage{textcomp}
\usepackage{gensymb}
\usepackage{adjustbox}
\usepackage{subcaption}
\usepackage{enumitem}
\usepackage{multicol}
\usepackage{listings}
\usepackage{url}
\usepackage{graphicx} % <-- needed for images
\def\UrlBreaks{\do\/\do-}

\usetheme{Boadilla}
\usecolortheme{lily}
\setbeamertemplate{footline}{
	\leavevmode%
	\hbox{%
		\begin{beamercolorbox}[wd=\paperwidth,ht=2ex,dp=1ex,right]{author in head/foot}%
			\insertframenumber{} / \inserttotalframenumber\hspace*{2ex}
	\end{beamercolorbox}}%
	\vskip0pt%
}
\setbeamertemplate{navigation symbols}{}

\lstset{
	frame=single,
	breaklines=true,
	columns=fullflexible,
	basicstyle=\ttfamily\tiny    % tiny font so code fits
}

\numberwithin{equation}{section}

% ---- your macros ----
\providecommand{\nCr}[2]{\,^{#1}C_{#2}}
\providecommand{\nPr}[2]{\,^{#1}P_{#2}}
\providecommand{\mbf}{\mathbf}
\providecommand{\pr}[1]{\ensuremath{\Pr\left(#1\right)}}
\providecommand{\qfunc}[1]{\ensuremath{Q\left(#1\right)}}
\providecommand{\sbrak}[1]{\ensuremath{{}\left[#1\right]}}
\providecommand{\lsbrak}[1]{\ensuremath{{}\left[#1\right.}}
\providecommand{\rsbrak}[1]{\ensuremath{\left.#1\right]}}
\providecommand{\brak}[1]{\ensuremath{\left(#1\right)}}
\providecommand{\lbrak}[1]{\ensuremath{\left(#1\right.}}
\providecommand{\rbrak}[1]{\ensuremath{\left.#1\right)}}
\providecommand{\cbrak}[1]{\ensuremath{\left\{#1\right\}}}
\providecommand{\lcbrak}[1]{\ensuremath{\left\{#1\right.}}
\providecommand{\rcbrak}[1]{\ensuremath{\left.#1\right\}}}
\theoremstyle{remark}
\newtheorem{rem}{Remark}
\newcommand{\sgn}{\mathop{\mathrm{sgn}}}
\providecommand{\abs}[1]{\left\vert#1\right\vert}
\providecommand{\res}[1]{\Res\displaylimits_{#1}}
\providecommand{\norm}[1]{\lVert#1\rVert}
\providecommand{\mtx}[1]{\mathbf{#1}}
\providecommand{\mean}[1]{E\left[ #1 \right]}
\providecommand{\fourier}{\overset{\mathcal{F}}{ \rightleftharpoons}}
\providecommand{\system}{\overset{\mathcal{H}}{ \longleftrightarrow}}
\providecommand{\dec}[2]{\ensuremath{\overset{#1}{\underset{#2}{\gtrless}}}}
\newcommand{\myvec}[1]{\ensuremath{\begin{pmatrix}#1\end{pmatrix}}}
\newcommand{\mydet}[1]{\ensuremath{\begin{vmatrix}#1\end{vmatrix}}}

\newenvironment{amatrix}[1]{%
	\left(\begin{array}{@{}*{#1}{c}|*{#1}{c}@{}}
	}{%
	\end{array}\right)
}

\newcommand{\myaugvec}[2]{\ensuremath{\begin{amatrix}{#1}#2\end{amatrix}}}
\let\vec\mathbf
% ---------------------

\title{12.55}
\author{ee25btech11042-Nipun Dasari}

\begin{document}
	
	\begin{frame}
		\titlepage
	\end{frame}
	
	\begin{frame}{Question}
		
		For a matrix \\
	\begin{align}
		\vec{M}=\begin{myvec}{\frac{4}{5}&\frac{3}{5}\\\frac{3}{5}&x}\end{myvec} \label{matrix}
	\end{align}
	The transpose of the matrix is equal to inverse of the matrix, i.e., 
	\begin{align}
		\vec{M}^\top = \vec{M}^{-1} \label{ortho}
	\end{align}
	The value of x is given by
		
	\end{frame}
	
	
	
	\begin{frame}{Solution}
		
			From \eqref{ortho} it is evident that the matrix \eqref{matrix} is an orthogonal matrix.\\
		We can proceed by the fact that columns of an orthogonal matrix are pairwise perpendicular
		\begin{align}
			\vec{M} = \begin{myvec}{\vec{c_1}&\vec{c_2}}\end{myvec}
		\end{align} 
		where
		\begin{align}
			\vec{c_1} = \begin{myvec}{\frac{4}{5}\\\frac{3}{5}}\end{myvec}\text{, }\vec{c_2} = \begin{myvec}{\frac{3}{5}\\x}\end{myvec}
		\end{align}
		
		
	\end{frame}
	
	\begin{frame}{Solution}
		
		
	By using \eqref{ortho}:
	\begin{align}
		\vec{c_1}^\top\vec{c_2} = 0\\
		\implies \begin{myvec}{\frac{4}{5}&\frac{3}{5}}\end{myvec}\begin{myvec}{\frac{3}{5}\\x}\end{myvec} = 0\\
		\implies \frac{12}{25} + \frac{3x}{5} = 0\\
		\implies x = -\frac{4}{5}
	\end{align}
		
	\end{frame}
	
	
		

	
\end{document}